\section{Introduction}
\label{sec:introduction}

Let $x_1,\dots,x_n$ be Boolean variables.  A \emph{proper signed
$k$-clause} is a disjunction of $k$ literals on $k$ distinct variables.
In the classical random $k$-SAT process, independent clauses are sampled
uniformly with replacement, and one studies the first prefix whose
conjunction is unsatisfiable.  Sharp threshold phenomena for random
$k$-SAT originate in the work of \citet{Friedgut1999}; quantitative lower
bounds on the width of the uniform transition were obtained by
\citet{Wilson2002}.  Recent work of \citet{Carenini2026Scaling} and
\citet{Carenini2026Polynomial} gives substantially sharper scaling-window
bounds and establishes the fixed-$k$ satisfiability conjecture.

The fluctuation problem for the first unsatisfiable prefix is finer than the
existence or width of a threshold.  In the uniform proper-clause model,
\citet{OpenAI2026KSATVariance} proved variance $\Theta_k(n)$ for fixed
$k\ge4$ and an $O(n\log n)$ upper bound for $k=3$.  The logarithmic loss in
the three-literal case was subsequently removed in
\citet{OpenAI2026ThreeSATVariance}, yielding variance $\Theta(n)$.  Their
upper-bound arguments combine coordinate omission, necessary tests obtained
by flipping variables of a deleted clause, and an occupation estimate for a
potential defined on the surviving assignment set.

The purpose of this paper is to show that the linear upper bound is stable
under bounded, time-dependent perturbations of the complete clause law.
The clause stream may be nonstationary, its variable and sign choices may
exhibit arbitrary dependence within a clause, and different times may use
different distributions.  Independence \emph{between} clause times is
retained.  The only structural assumption is a two-sided pointwise density
bound relative to the uniform proper-clause law.

\subsection{Model and main result}

Let $\cC_{n,k}$ denote the set of proper signed $k$-clauses on
$x_1,\dots,x_n$, and let $U_{n,k}$ be the uniform probability measure on
$\cC_{n,k}$:
\[
    U_{n,k}(C)=\frac{1}{2^k\binom nk}
    \qquad(C\in\cC_{n,k}).
\]
For every $n$, let $C_1,C_2,\dots$ be independent clauses, where $C_j$ has
law $\mu_{n,j}$ on $\cC_{n,k}$.  We assume that fixed constants
\[
    0<a\le 1\le b<\infty
\]
satisfy
\begin{equation}
    a\,U_{n,k}(C)
    \le \mu_{n,j}(C)
    \le b\,U_{n,k}(C)
    \label{eq:uniform-comparability}
\end{equation}
for every $n$, every time $j$, and every $C\in\cC_{n,k}$.  Put
\[
    F_m=\bigwedge_{j=1}^m C_j,
    \qquad
    H_n=\min\{m\ge1:F_m\text{ is unsatisfiable}\}.
\]

\begin{theorem}[Robust linear variance bound]
\label{thm:main}
Fix $k\ge3$ and constants $0<a\le1\le b<\infty$.  Under
\eqref{eq:uniform-comparability}, there is a constant
$K=K(a,b,k)<\infty$ such that
\[
    \Var(H_n)\le Kn
    \qquad(n\ge k).
\]
More generally, for every fixed $\rho<\infty$, there is a constant
$K_\rho=K_\rho(a,b,k)<\infty$ such that
\[
    \Var(H_n\wedge L)\le K_\rho n
\]
for all $n\ge k$ and all integers $1\le L\le \rho n$.
\end{theorem}

The constants are uniform over the entire triangular array
$(\mu_{n,j})_{n,j}$.  In particular, deterministic time dependence is
allowed without any regularity assumption in $j$.

The proof yields a more local theorem.  For a $k$-set $R$ of roots, reveal
from a clause the subset of $R$ that it uses and the literal signs on that
subset.  Conditional on this revealed root trace, the remaining literals are
called the hidden nonroot part.

\begin{theorem}[Local robustness criterion]\label{thm:local-criterion}
Fix $k\ge3$.  Suppose the clauses are independent across time and that there
are constants
\[
 B_{\mathrm{o}},B_{\mathrm{t}}\ge1,\qquad
 \gamma_1,\gamma_2,\delta>0
\]
with the following properties, uniformly in $n$, the time index, every
$k$-set $R\subseteq[n]$, and every relevant root trace of positive
probability:
\begin{enumerate}[label=\textnormal{(\roman*)}]
\item conditional on avoiding $R$, the hidden proper $k$-clause on the
      $n-k$ nonroots has density at least $B_{\mathrm{o}}^{-1}$ relative to the
      uniform law;
\item conditional on using exactly one specified root with a specified sign,
      the hidden proper $(k-1)$-clause has density at most $B_{\mathrm{t}}$
      relative to the uniform law;
\item a clause meets $R$ with probability at most $\gamma_1/n$, and meets at
      least two elements of $R$ with probability at most
      $\gamma_2/[n(n-1)]$;
\item for every fixed assignment and every time, the sampled clause falsifies
      that assignment with probability at least $\delta$.
\end{enumerate}
Then, for every fixed $\rho<\infty$,
\[
 \Var(H_n\wedge L)=O(n)\qquad(1\le L\le\rho n),
\]
where the constant depends only on
$k,\rho,B_{\mathrm{o}},B_{\mathrm{t}},\gamma_1,\gamma_2$.  Moreover
$\Var(H_n)=O(n)$, with dependence also on $\delta$.
\end{theorem}

Under \eqref{eq:uniform-comparability}, the theorem applies with
\begin{equation}\label{eq:local-parameters}
 B_{\mathrm{o}}=B_{\mathrm{t}}=\frac ba,
 \qquad \gamma_1=bk^2,
 \qquad \gamma_2=\frac{bk^2(k-1)^2}{2},
 \qquad \delta=a2^{-k}.
\end{equation}
Indeed, conditional densities lie in $[a/b,b/a]$; the incidence estimates
follow by union bounds over roots and root pairs; and a fixed assignment
falsifies a uniform proper clause with probability $2^{-k}$.  Thus
Theorem~\ref{thm:main} is a direct corollary of the local criterion.

The variance estimate immediately controls every fixed central quantile
window.  Define
\[
    s_n(m)=\Prb(H_n>m),
    \qquad
    r_n(p)=\min\{m\ge0:s_n(m)\le p\},
    \quad 0<p<1.
\]

\begin{corollary}[Central windows]
\label{cor:central-window}
Under the hypotheses of Theorem~\ref{thm:main}, for every fixed
$0<\varepsilon<1/2$,
\[
    r_n(\varepsilon)-r_n(1-\varepsilon)
    =O_{a,b,k,\varepsilon}(\sqrt n).
\]
\end{corollary}

We prove only an upper variance bound.  The linear lower bound known for the
uniform model uses additional information about separation of satisfiability
levels and does not automatically persist for every nonidentical stream
satisfying \eqref{eq:uniform-comparability}.  Although the lower density bound
permits the mixture representation
\[
 \mu_{n,j}=aU_{n,k}+(1-a)\nu_{n,j},
\]
Bernoulli thinning alone does not give a variance lower bound: the residual
clauses may determine the hitting time before fluctuations of the uniform
component become pivotal.  Establishing a matching lower bound therefore
requires a separate anti-concentration or level-separation argument and
remains open in this generality.

\subsection{Proof mechanism}

Fix one deleted clause and call its $k$ variables \emph{roots}.  If that
clause is pivotal, then every surviving solution has the root assignment
that falsifies it.  Flipping an available root produces a necessary family
of $(k-1)$-clause tests on the nonroot solution set.  If $p$ denotes the
probability that $d$ independent uniform $(k-1)$-clauses kill that set, the
$p$-dependence of the argument is
\[
    \underbrace{p^k}_{\text{necessary root tests}}
    \times
    \underbrace{p^{-k/(k-1)}}_{\text{remaining duration}}
    \le
    \underbrace{p^{k/(k-1)}}_{\text{controlled occupation}}.
\]
The inequality is exactly the exponent condition
\[
    k-\frac{k}{k-1}\ge \frac{k}{k-1},
\]
which holds for every $k\ge3$ and is an equality at $k=3$.

The local hypotheses have distinct roles.  The lower conditional density
$B_{\mathrm{o}}^{-1}$ drives the occupation potential at ordinary clauses.  The
upper conditional test density $B_{\mathrm{t}}$ costs a factor at most
$B_{\mathrm{t}}^r$ for a family of $r$ hidden root tests.  The incidence and
collision bounds make the total root-incidence count exponentially
integrable and collisions rare, while $\delta$ supplies the tail estimate
used to remove the cap.  The two-sided pointwise hypothesis is therefore a
convenient sufficient condition, not the intrinsic formulation of the
argument.

Section~\ref{sec:potential} proves the general-$k$ killing-potential
inequality.  Sections~\ref{sec:omission} and \ref{sec:exposure} set up the
clause deletion and conditional exposure.  Sections~\ref{sec:tests} and
\ref{sec:occupation} prove the root-test and occupation estimates.
Section~\ref{sec:delay} obtains a uniform second-moment bound for each
deletion delay, and Section~\ref{sec:uncapping} removes the linear cap.
